\documentclass{article}
\usepackage{hyperref}
\usepackage{Style}

\nocite{*} % Comentar si quiero citar
%\addbibresource{bibliografia.bib} % Quitar el comentado si quiero usar bibliografia

\begin{document}

\begin{minipage}{2.5cm}
    \includegraphics[width=2cm]{imagen_puc.jpg}
\end{minipage}
\begin{minipage}{12cm}
    {\sc Pontificia Universidad Católica de Chile\\
    Facultad de Matemáticas\\
    Departamento de Matemática\\
    Profesor: Manuel Cabezas -- Ayudante: Benjamín Mateluna}
\end{minipage}
\vspace{2ex}

{\centerline{\bf Introducción a la Geometría - MAT1304}
\centerline{\bf Ayudantía 18}}
\centerline{\bf 14 de mayo de 2026}

\vspace{1cm}
\noindent\textbf{Problema 1.} Halle un polinomio $p(x)\in\R[x]$ de grado 5, con raíces $-1+i,
\sqrt{2}i$ y $1$ tal que $p(-1)=-6$.

\vspace{5mm}
\noindent\textbf{Problema 2.} Resuelva las siguientes ecuaciones
\begin{enumerate}
    \item $z^{5}=\sqrt{3}+i$.
    \item $z^{7}-2iz^{4}-iz^{3}-2=0$.
\end{enumerate}

\vspace{5mm}
\noindent\textbf{Problema 3.} Sea $p(x)=x^{2}-a$ con $a\in\Z$, muestre que $p(x)$ es reducible en 
$\Z[x]$ si y solo si existe $\alpha\in\Z$ tal que $\alpha^{2}=a$.

\vspace{5mm}
\begin{dfn}
    Sea $F\in\C[x,y]$ \footnote{
    El conjunto $\C[x,y]$ corresponde a los polinomios en dos variables con coeficientes en $\C$, 
    estos se suman y multiplican de igual manera que los polinomios en una variable.
}, se dice homogéneo si para todo $\lambda\in\C$, existes $n\in\N$ tal que
    \begin{equation*}
        F(\lambda x,\lambda y)=\lambda^{n}F(x,y)
    \end{equation*}
\end{dfn}

\noindent\textbf{Problema 4.} Demuestre que todo polinomio homogéneo en $\C[x,y]$ se puede 
escribir como el producto de polinomios lineales homogéneos.

\vspace{1cm}
\noindent\textbf{\underline{Ejercicios Propuestos:}}

\vspace{5mm}
\noindent\textbf{Extra 1.} Muestre que $p(x)$ y $q(x)$ tienen un factor común si y solo si tienen
una raíz compleja en común.

\vspace{5mm}
\noindent\textbf{Extra 2.} Factorizar $p(x,y)=x^{3}-2x^{2}y-y^{2}x+2y^{3}$ en $\C[x,y]$.

%\printbibliography % Quitar el comentado si quiero usar bibliografia

\end{document}
